% ============================================
% Postdoctoral Associate Cover Letter – University at Albany (SUNY)
% ============================================

\documentclass[11pt,letterpaper]{letter}

\usepackage{geometry}
\usepackage{microtype}
\usepackage[hidelinks]{hyperref}

\geometry{left=25mm,right=25mm,top=25mm,bottom=25mm}

\signature{Decory Edwards}
\address{
Department of Economics \\
Johns Hopkins University \\
Baltimore, MD 21218 \\
\texttt{dedwar65@jh.edu}
}

\begin{document}

\begin{letter}{Search Committee \\
Department of Economics \\
University at Albany, SUNY \\
Albany, NY}

\opening{Dear Members of the Search Committee,}

I am writing to apply for the Postdoctoral Associate position in Economics at the University at Albany, SUNY. I am a Ph.D.\ candidate in Economics at Johns Hopkins University, advised by Professors Christopher D.~Carroll, Nicholas W.~Papageorge, and Stelios Fourakis, and I expect to receive my degree in May~2026. My research fields are macroeconomics, wealth and income dynamics, and computational economics.

My research agenda focuses on understanding the determinants of wealth inequality and the mechanisms through which heterogeneity in returns to wealth shapes aggregate outcomes. In my job-market paper, I estimate the degree of heterogeneity in individual portfolio returns using micro-data from the Survey of Consumer Finances and simulate a structural model in which households face idiosyncratic return risk. I show that allowing for persistent heterogeneity in returns substantially improves the model’s ability to match the empirical distribution of wealth, offering a tractable way to connect micro-level return dispersion to macroeconomic inequality.

In a second project, I use the Health and Retirement Study to examine the relationship between interpersonal trust and portfolio performance. Building on the literature that finds a hump-shaped relationship between trust and income, I show that a similar relationship holds for individual returns to wealth measured in the HRS data, highlighting a behavioral channel through which social attitudes shape saving and investment outcomes. This work also provides new estimates of the persistent component of returns to net worth, which motivate the calibration of heterogeneity in my structural model.

The postdoctoral position at the University at Albany is particularly appealing because of its emphasis on active research, collaboration, and the development of early-career scholars. The department’s strengths in applied microeconomics, macroeconomics, and quantitative methods provide an excellent environment for advancing my research agenda and for deepening connections between micro-level evidence and aggregate economic dynamics. I would value the opportunity to engage with faculty whose work spans empirical analysis, policy relevance, and methodological rigor, and to contribute to seminars and collaborative projects during the appointment.

As a postdoctoral associate, I would focus on extending my current research pipeline toward publication, developing new projects that build on my structural and empirical work, and participating actively in the intellectual life of the department. I am especially interested in using the postdoctoral period to sharpen the policy implications of my research and to broaden its relevance across macroeconomic and applied settings.

Thank you very much for your time and consideration. I would welcome the opportunity to discuss my application further.

\closing{Sincerely,}

\end{letter}
\end{document}
